\section{GitLab}
\label{sec:c51-gitlab}

GitLab là nền tảng DevOps tích hợp toàn diện, kết hợp quản lý mã nguồn Git,
theo dõi công việc, review code, kiểm thử và triển khai trong một ứng dụng duy nhất.
Với đồ án ASPX trên .NET Framework 4.8.1, GitLab giúp nhóm ba sinh viên làm việc
trên cùng một kho mã mà vẫn kiểm soát được chất lượng qua Merge Request,
Approval rules, protected branch và pipeline tự động.
Chương này trình bày các khái niệm GitLab cơ bản nhất mà nhóm đã dùng thật cho đồ án:
Groups, Projects, Issues, Boards, Milestones, Wiki, Snippets, Merge Request và bảo vệ nhánh.
Nội dung phần CI/CD nâng cao với Windows runner và \texttt{msbuild} được trình bày ở chương tiếp theo,
để chương này tập trung vào quy trình cộng tác hằng ngày \cite{gitlabdocs}.

\subsection{Tổng quan GitLab: nền tảng DevOps tích hợp}
\label{subsec:c51-tong-quan}

GitLab đi theo triết lý DevOps Platform: một vòng đời khép kín từ ý tưởng tới vận hành.
Thay vì ghép nhiều công cụ rời rạc, nhóm chỉ cần một nơi duy nhất để lập kế hoạch,
viết code, review, kiểm thử, phát hành và giám sát.
Cách tiếp cận này đặc biệt phù hợp với đồ án môn học, vì nhóm nhỏ, thời gian ngắn,
nhưng vẫn phải thể hiện được quy trình chuyên nghiệp giống doanh nghiệp \cite{gitlabdocs}.

Ba trụ cột mà nhóm khai thác nhiều nhất gồm quản lý mã nguồn Git, cộng tác qua Merge Request,
và quản trị công việc qua Issues cùng Boards.
Mọi thay đổi trên form \texttt{Default.aspx}, code-behind \texttt{Default.aspx.cs} hay
\texttt{Web.config} đều đi qua một Merge Request có người review, thay vì đẩy thẳng lên nhánh chính.
Nhờ vậy, lịch sử commit sạch, dễ truy vết ai sửa gì, vì sao sửa, và có thể hoàn tác an toàn.

GitLab được cung cấp theo hai mô hình triển khai chính.
Mô hình SaaS tại \url{https://gitlab.com} do GitLab vận hành sẵn, nhóm chỉ cần đăng ký tài khoản là dùng ngay.
Mô hình self-hosted cho phép khoa hoặc doanh nghiệp tự cài đặt GitLab lên máy chủ riêng để toàn quyền kiểm soát.
Bảng \ref{tab:c51-saas-selfhosted} so sánh hai mô hình theo năm tiêu chí quan trọng với đồ án môn học \cite{gitlabdocs}.

\begin{table}[H]
\centering
\caption{So sánh GitLab SaaS và GitLab self-hosted}
\label{tab:c51-saas-selfhosted}
\begin{tabular}{lll}
\toprule
Tiêu chí & SaaS gitlab.com & Self-hosted \\
\midrule
Cài đặt và vận hành & Đăng ký là dùng ngay & Tự cài đặt và bảo trì máy chủ \\
Chi phí & Miễn phí cho nhóm nhỏ & Tốn máy chủ và công quản trị \\
Kiểm soát dữ liệu & Dữ liệu đặt trên cloud & Dữ liệu nằm hoàn toàn nội bộ \\
Runner cho ASPX 4.8.1 & Dùng runner cloud có sẵn & Tự gắn Windows runner riêng \\
Nâng cấp và tùy biến & Tự động nâng cấp & Chủ động chọn phiên bản \\
\bottomrule
\end{tabular}
\end{table}

Với đồ án ASPX của nhóm, mô hình SaaS là lựa chọn hợp lý vì không tốn máy chủ,
vẫn đủ tính năng cho ba thành viên, và dễ chia sẻ link công khai cho giảng viên chấm bài.
Nhóm tạo một Group chung, bên trong là Project chứa toàn bộ source Web Forms,
bật Issues và Merge Request ngay từ đầu để mọi trao đổi đều lưu lại trên GitLab.
Lệnh dưới đây cho thấy thao tác đầu tiên của thành viên mới: clone kho về máy,
tạo nhánh tính năng và đẩy lên GitLab để chuẩn bị mở Merge Request.

\begin{lstlisting}[language=bash, caption={Clone kho ASPX và đẩy nhánh tính năng lên GitLab}, label={lst:c51-clone-push}]
# Clone the ASPX demo project from GitLab
git clone https://gitlab.com/nhom3/aspxdemo-net481.git
cd aspxdemo-net481
git config user.name "Sinh Vien 3"
git config user.email "sv3@example.com"
git checkout -b feature/login-aspx
# Edit Default.aspx and Default.aspx.cs, then commit
git add Default.aspx Default.aspx.cs Web.config
git commit -m "feat: them form dang nhap ASPX"
git push -u origin feature/login-aspx
\end{lstlisting}

Sau khi đẩy nhánh, GitLab hiển thị nút tạo Merge Request ngay trên trang Project.
Thành viên chỉ cần bấm nút, ghi mô tả thay đổi, gán reviewer là một bạn khác trong nhóm,
rồi chờ pipeline chạy và chờ approve trước khi merge.
Quy trình này tuy thêm một bước so với đẩy thẳng, nhưng loại bỏ phần lớn lỗi do sửa nhầm file dùng chung.

\subsection{Groups, Projects, Issues, Boards, Milestones, Wiki, Snippets}
\label{subsec:c51-cong-cu-cong-tac}

Bảy khái niệm dưới đây tạo thành bộ khung cộng tác hằng ngày của nhóm trên GitLab.
Mỗi khái niệm được giải thích gắn trực tiếp với đồ án ASPX, để người đọc hình dung ngay cách dùng.

\begin{description}
\item[Groups] Group là đơn vị tổ chức cao nhất mà nhóm dùng, đại diện cho cả nhóm ba sinh viên.
Nhóm đặt tên Group theo tên đề tài, thêm cả ba thành viên với vai trò Developer,
còn giảng viên được thêm với vai trò Reporter để xem code và pipeline mà không sửa được.
Mọi Project con trong Group đều kế thừa thành viên, giúp phân quyền một lần cho tất cả kho.

\item[Projects] Project là một kho Git kèm đầy đủ Issues, Merge Request, Wiki và pipeline.
Đồ án ASPX chỉ cần một Project duy nhất chứa toàn bộ solution .NET Framework 4.8.1,
gồm thư mục web, file \texttt{Web.config}, file solution và thư mục tài liệu.
Nhóm bật tính năng Issues, Wiki và Container Registry ngay khi tạo Project,
để không phải cấu hình rời rạc về sau.

\item[Issues] Issue là đơn vị công việc nhỏ nhất, dùng để ghi nhận việc cần làm hoặc lỗi cần sửa.
Ví dụ, khi form đăng nhập báo lỗi sai chuỗi kết nối, nhóm tạo Issue với tiêu đề mô tả lỗi,
gán nhãn \texttt{bug} và \texttt{priority-high}, rồi gán cho thành viên phụ trách xử lý.
Mỗi Issue đều có số thứ tự, người được gán, nhãn, hạn chót và luồng thảo luận riêng.
Nhờ vậy, không còn tình trạng giao việc bằng tin nhắn rồi quên mất ai làm gì.

\item[Boards] Board là bảng Kanban trực quan, kéo thả các Issue qua các cột như To do, Doing và Done.
Nhóm họp nhanh mỗi tối, mở Board ra là thấy ngay ai đang làm dở việc gì, việc nào tắc.
Khi bắt đầu sửa form đăng nhập, thành viên kéo Issue sang cột Doing và tự gán cho mình.
Khi Merge Request tương ứng được merge, Issue tự đóng và thẻ biến mất khỏi Board.
Cách làm này giúp tiến độ đồ án luôn hiển thị công khai, không cần báo cáo thủ công.

\item[Milestones] Milestone là mốc thời gian gom nhiều Issue lại, tương ứng với một đợt nộp bài.
Nhóm đặt ba mốc: bản nháp giao diện, bản tích hợp đầy đủ chức năng, và bản nộp cuối cùng.
Mỗi Issue khi tạo đều được gắn vào một mốc, nên trang Milestone luôn cho biết còn bao nhiêu việc chưa xong.
Giảng viên chỉ cần mở trang Milestone là đánh giá được tiến độ, thay vì phải hỏi từng thành viên.

\item[Wiki] Wiki là nơi viết tài liệu nội bộ: hướng dẫn cài môi trường, cách cấu hình IIS,
cách khôi phục cơ sở dữ liệu mẫu và danh sách tài khoản thử nghiệm.
Khác với file văn bản trong kho code, Wiki có lịch sử sửa đổi riêng và sửa trực tiếp trên trình duyệt.
Nhóm ghi vào Wiki các bước build bằng \texttt{msbuild} cho .NET Framework 4.8.1,
để thành viên mới chỉ cần đọc Wiki là dựng được môi trường trong một buổi.

\item[Snippets] Snippet là đoạn code hoặc cấu hình ngắn được chia sẻ nhanh trong nhóm.
Ví dụ, đoạn chuỗi kết nối mẫu, đoạn cấu hình lại \texttt{Web.config} hay hàm kiểm tra đăng nhập viết bằng C\#,
đều được lưu thành Snippet thay vì gửi qua tin nhắn dễ thất lạc.
Mỗi Snippet có tiêu đề, định dạng ngôn ngữ và lịch sử phiên bản, dùng lại được nhiều lần.
\end{description}

Khi đã quen, nhóm tạo Issue trực tiếp bằng GitLab API trong các script tự động.
Ví dụ dưới đây tạo một Issue báo lỗi xung đột \texttt{Web.config} bằng lệnh \verb|curl|,
rất hữu ích khi pipeline phát hiện lỗi và muốn ghi nhận tự động.

\begin{lstlisting}[language=bash, caption={Tạo Issue tự động bằng GitLab API}, label={lst:c51-create-issue}]
# Create an issue via GitLab API (needs a personal token)
export GITLAB_TOKEN="glpat-xxxxxxxx"
export PROJECT_ID="12345678"
curl --request POST \
  --header "PRIVATE-TOKEN: $GITLAB_TOKEN" \
  --data "title=Fix Web.config conflict on develop" \
  --data "description=Merge develop gay xung dot muc connectionStrings" \
  --data "labels=bug,ASPX,priority-high" \
  "https://gitlab.com/api/v4/projects/$PROJECT_ID/issues"
\end{lstlisting}

\subsection{Merge Request và Approval rules}
\label{subsec:c51-mr-approval}

Merge Request, viết tắt là MR, là cơ chế review code trung tâm của GitLab,
tương đương với Pull Request trên GitHub nhưng gắn chặt hơn với pipeline và Approval rules.
Mọi thay đổi cho đồ án ASPX, từ sửa giao diện tới đổi chuỗi kết nối, đều phải đi qua MR.
Quy định này giúp nhóm ba người không bao giờ ghi đè lẫn nhau trên nhánh chính.

Vòng đời điển hình của một MR trong đồ án gồm sáu bước.

\begin{enumerate}
\item Tạo nhánh tính năng từ \texttt{develop}, đặt tên rõ ràng như \texttt{feature-login}.
\item Viết code, commit theo Conventional Commits rồi đẩy nhánh lên GitLab.
\item Mở MR từ nhánh tính năng về \texttt{develop} hoặc \texttt{main}, ghi mô tả và ảnh chụp màn hình.
\item Pipeline tự động chạy build và kiểm tra cơ bản trên chính MR đó.
\item Một hoặc hai thành viên khác review code, bình luận và approve.
\item Merge MR khi đủ approve và pipeline báo thành công, sau đó xóa nhánh tính năng.
\end{enumerate}

Approval rules là quy tắc bắt buộc về số người phê duyệt trước khi được merge.
Với đồ án môn học, nhóm cấu hình nhánh \texttt{main} yêu cầu hai approve,
nghĩa là cả hai thành viên còn lại đều phải đồng ý thì MR mới được merge.
Quy tắc này ngăn trường hợp một người tự merge code chưa ai xem vào bản nộp cuối.
Nhánh \texttt{develop} yêu cầu nhẹ hơn, chỉ cần một approve để tốc độ làm bài không bị chậm.

Hình \ref{fig:c51-mr-flow} tóm tắt luồng MR từ nhánh tính năng tới nhánh chính,
nhấn mạnh hai chốt chặn quan trọng là pipeline và số lượng approve.

\begin{figure}[H]
\centering
\begin{tikzpicture}[>=Stealth, thick, box/.style={draw, rounded corners, fill=blue!10, minimum width=2.4cm, minimum height=0.9cm, align=center, font=\small}]
\node[box] (feat) {Nhanh tinh nang};
\node[box, right=1.2cm of feat] (mr) {Mo MR};
\node[box, right=1.2cm of mr] (review) {Review va\\Approve};
\node[box, right=1.2cm of review] (merge) {Merge vao\\main};
\draw[->] (feat) -- node[midway, above, font=\footnotesize] {push} (mr);
\draw[->] (mr) -- node[midway, above, font=\footnotesize] {pipeline} (review);
\draw[->] (review) -- node[midway, above, font=\footnotesize] {2 approve} (merge);
\end{tikzpicture}
\caption{Luong Merge Request tu nhanh tinh nang toi nhanh chinh}
\label{fig:c51-mr-flow}
\end{figure}

Ví dụ dưới đây tạo MR bằng API rồi tạo Approval rule yêu cầu hai người phê duyệt.
Trong thực tế nhóm cấu hình bằng giao diện web, nhưng các lệnh này giúp hiểu rõ bản chất
và có thể dùng lại khi dựng Project mới cho môn học khác \cite{gitlabdocs}.

\begin{lstlisting}[language=bash, caption={Tạo Merge Request và Approval rule bằng API}, label={lst:c51-mr-approve}]
# Create a merge request from feature branch to main
curl --request POST \
  --header "PRIVATE-TOKEN: $GITLAB_TOKEN" \
  --data "source_branch=feature/login-aspx" \
  --data "target_branch=main" \
  --data "title=Merge login form to main" \
  "https://gitlab.com/api/v4/projects/$PROJECT_ID/merge_requests"
# Require 2 approvals before merge (approval rule)
curl --request POST \
  --header "PRIVATE-TOKEN: $GITLAB_TOKEN" \
  --data "name=Require 2 approvers" \
  --data "approvals_required=2" \
  "https://gitlab.com/api/v4/projects/$PROJECT_ID/approval_rules"
\end{lstlisting}

\subsection{Bảo vệ nhánh và tag}
\label{subsec:c51-bao-ve-nhanh}

Bảo vệ nhánh là lớp khóa cuối cùng giữ cho các nhánh quan trọng không bị ghi đè trực tiếp.
GitLab phân biệt hai quyền riêng: Allowed to push quyết định ai được đẩy trực tiếp lên nhánh,
còn Allowed to merge quyết định ai được bấm nút merge MR vào nhánh đó.
Tách bạch hai quyền này giúp nhóm áp dụng quy tắc nghiêm cho bản nộp cuối,
nhưng vẫn linh hoạt cho nhánh phát triển hằng ngày \cite{gitlabdocs}.

Nhóm áp dụng cấu hình trong Bảng \ref{tab:c51-protected} cho hai nhánh quan trọng nhất.

\begin{table}[H]
\centering
\caption{Cấu hình bảo vệ cho hai nhánh main và develop}
\label{tab:c51-protected}
\begin{tabular}{llll}
\toprule
Nhánh & Allowed to push & Allowed to merge & Yêu cầu bổ sung \\
\midrule
main & Không ai & Maintainers & 2 approve, pipeline thành công \\
develop & Developers & Developers & 1 approve, pipeline thành công \\
\bottomrule
\end{tabular}
\end{table}

Với cấu hình trên, nhánh \texttt{main} không cho phép đẩy trực tiếp bằng lệnh \verb|git push|,
mọi thay đổi bắt buộc phải qua MR và được hai người phê duyệt.
Nhánh \texttt{develop} cho phép thành viên đẩy trực tiếp khi sửa lỗi nhỏ,
nhưng các tính năng lớn như thêm trang quản trị vẫn phải mở MR để có người kiểm tra chéo.
Cách phân quyền này cân bằng giữa an toàn và tốc độ, rất phù hợp với nhóm ba người.

Ngoài nhánh, GitLab còn cho phép bảo vệ tag để khóa các mốc phát hành.
Nhóm gắn tag \texttt{v1.0} cho bản nộp giữa kỳ và tag \texttt{v2.0} cho bản nộp cuối,
sau đó cấu hình chỉ Maintainers mới được tạo tag, tránh việc gắn nhầm tag lên commit lỗi.
Khi cần chấm lại bản cũ, giảng viên chỉ cần checkout đúng tag là có đúng bản đã nộp.

Các lệnh dưới đây minh họa cách bảo vệ nhánh bằng API, tương đương với thao tác trên giao diện web.
Mức quyền số 0 nghĩa là không ai được đẩy trực tiếp, mức 30 là Developers, mức 40 là Maintainers.

\begin{lstlisting}[language=bash, caption={Bảo vệ nhánh main và develop bằng API}, label={lst:c51-protect}]
# Protect main: no direct push, only maintainers can merge
curl --request POST \
  --header "PRIVATE-TOKEN: $GITLAB_TOKEN" \
  --data "name=main" \
  --data "push_access_level=0" \
  --data "merge_access_level=40" \
  "https://gitlab.com/api/v4/projects/$PROJECT_ID/protected_branches"
# Protect develop: developers can push and merge with review
curl --request POST \
  --header "PRIVATE-TOKEN: $GITLAB_TOKEN" \
  --data "name=develop" \
  --data "push_access_level=30" \
  --data "merge_access_level=30" \
  "https://gitlab.com/api/v4/projects/$PROJECT_ID/protected_branches"
\end{lstlisting}

Tóm lại, GitLab cung cấp cho nhóm đầy đủ công cụ cộng tác từ Group tới Snippets,
với Merge Request và Approval rules làm trung tâm kiểm soát chất lượng.
Khi kết hợp thêm bảo vệ nhánh và tag, nhóm ba sinh viên có thể làm việc song song
mà không sợ mất code hay ghi đè bản nộp cuối của nhau.
Đây chính là nền tảng để chương tiếp theo xây dựng pipeline CI/CD tự động cho đồ án ASPX 4.8.1.
